\sezione{9}{Esempio 2: pipeline, validazione e tuning}{Parte III · il flusso di lavoro «professionale»}

\begin{inbreve}
Stesso dataset dell'esempio 1, ma con il flusso di lavoro completo del capitolo 7: \hl{Pipeline} con preprocessing, \hl{validazione incrociata stratificata}, \hl{ricerca a griglia} degli iperparametri, valutazione con \key{AUC}, \hl{importanza per permutazione} e \hl{salvataggio} del modello. È lo scheletro da copiare per i tuoi progetti.
\end{inbreve}

\subsection{Il codice}

\codicefile{es2_pipeline_tuning.py}{1}{47}

\subsection{Spiegazione}

\begin{tecnica}[Pipeline \ \riga{22--25}]
La \texttt{Pipeline} è una lista di coppie \texttt{(nome, oggetto)}. L'ultimo elemento è il modello; i precedenti sono trasformazioni. Qui \texttt{SimpleImputer(strategy="median")} sostituisce eventuali valori mancanti con la mediana della colonna, calcolata \key{solo sui dati di training} di ogni fold. Questo dataset non ha valori mancanti: l'imputer è inserito per mostrare dove va il preprocessing in un progetto reale.
\end{tecnica}

\begin{tecnica}[Validazione incrociata \ \riga{28--31}]
\texttt{StratifiedKFold(n\_splits=5, shuffle=True)} divide il training in 5 parti con la stessa proporzione di classi. \texttt{cross\_val\_score} addestra 5 volte la pipeline (su 4 parti) e la valuta sulla parte esclusa. Si usa l'\key{AUC} (area sotto la curva ROC): misura quanto bene il modello \emph{ordina} i pazienti per probabilità, indipendentemente dalla soglia; vale 0.5 per un modello casuale e 1 per uno perfetto.
\end{tecnica}

\begin{tecnica}[Ricerca a griglia \ \riga{34--43}]
\texttt{griglia} elenca i valori da provare. Le chiavi hanno la forma \texttt{passo\_\_parametro}: \texttt{rf\_\_max\_features} è il parametro \texttt{max\_features} del passo chiamato \texttt{"rf"}. \texttt{GridSearchCV} prova tutte le $3\times3\times2=18$ combinazioni, ognuna con 5 fold: in totale 90 addestramenti, più uno finale. Alla fine riaddestra automaticamente la combinazione migliore su \emph{tutto} il training (\texttt{refit=True} di default) e la mette in \texttt{best\_estimator\_}.
\end{tecnica}

\begin{tecnica}[Valutazione finale \ \riga{45--47}]
Solo ora si usa il test set, una volta. \texttt{predict\_proba(X\_te)[:, 1]} prende la colonna della classe 1 (benigno): \texttt{roc\_auc\_score} vuole le probabilità della classe positiva.
\end{tecnica}

\begin{terminale}{python es2_pipeline_tuning.py  (prima parte)}
AUC per fold: [0.9884 0.99   0.9679 0.9947 0.9991]
AUC media   : ¤0.9880 +/- 0.0107¤
Migliori parametri: {'rf__max_depth': None, 'rf__max_features': 'sqrt', 'rf__min_samples_leaf': 3}
AUC CV migliore   : 0.9887
Combinazioni provate: 18
AUC sul test      : ¤0.9931¤
\end{terminale}

\begin{esempio}[Cosa ci dice l'output]
\begin{itemize}
  \item La variabilità tra fold (da 0.968 a 0.999) è più grande del guadagno del tuning (da 0.9880 a 0.9887): \hl{il tuning qui conta poco}, il Random Forest con i default è già vicino al meglio. È una situazione frequente.
  \item La combinazione scelta è $m=\sqrt p$, alberi non limitati in profondità e \texttt{min\_samples\_leaf=3}: foglie un po' più grandi rendono il modello leggermente più liscio.
  \item AUC sul test 0.993: coerente con la stima CV, nessun segnale di overfitting nella scelta degli iperparametri.
\end{itemize}
\end{esempio}

\subsection{Importanze e salvataggio}

\codicefile{es2_pipeline_tuning.py}{49}{61}

\begin{tecnica}[Spiegazione]
\begin{itemize}
  \item \riga{50} \texttt{named\_steps["rf"]} estrae il Random Forest dalla pipeline, per leggerne \texttt{feature\_importances\_} (\riga{51}). Con \texttt{pd.Series(..., index=X.columns)} ogni importanza è associata al nome della colonna.
  \item \riga{52--53} \texttt{permutation\_importance} si applica alla \key{pipeline intera} sul \key{test set}, con 20 permutazioni per colonna e la stessa metrica (AUC).
  \item \riga{59--61} \texttt{joblib.dump} salva su disco tutta la pipeline (imputer + foresta); \texttt{joblib.load} la ricarica. Il controllo finale verifica che le previsioni siano identiche.
\end{itemize}
\end{tecnica}

\begin{terminale}{python es2_pipeline_tuning.py  (seconda parte)}
Top 5 MDI:
 worst perimeter         0.1530
worst area              0.1504
worst concave points    0.1172
mean concave points     0.1046
worst radius            0.0761
Top 5 permutazione:
 worst area              0.0043
worst perimeter         0.0038
worst concave points    0.0023
mean concave points     0.0012
worst texture           0.0005
Stesse predizioni: True
\end{terminale}

Il confronto grafico tra le due importanze è nella figura del capitolo 6: le prime feature coincidono, ma la permutazione dà valori molto più piccoli perché le feature sono correlate.

\begin{attenzione}[Sicurezza dei file joblib]
Un file \texttt{.joblib} (come un \texttt{.pkl}) può eseguire codice arbitrario quando viene caricato: \hl{carica solo modelli che hai creato tu} o che provengono da fonti fidate. Inoltre un modello salvato con una versione di scikit-learn potrebbe non caricarsi con un'altra: annota la versione (ad esempio in \texttt{requirements.txt}).
\end{attenzione}

\begin{consiglio}[GridSearch o RandomizedSearch?]
Con pochi parametri e pochi valori va bene \texttt{GridSearchCV}. Quando le combinazioni sono tante, \texttt{RandomizedSearchCV} con \texttt{n\_iter} fissato (come nel capitolo 7) esplora lo stesso spazio con un costo controllato e spesso trova risultati altrettanto buoni.
\end{consiglio}

\begin{appunti}{Riepilogo del capitolo}
\begin{itemize}
  \item \texttt{Pipeline([("imputer", ...), ("rf", ...)])}: preprocessing dentro, niente data leakage.
  \item \texttt{StratifiedKFold} + \texttt{cross\_val\_score} per stimare la prestazione con incertezza.
  \item \texttt{GridSearchCV(pipe, griglia, cv=cv, scoring=...)}; chiavi \texttt{passo\_\_parametro}; risultato in \texttt{best\_estimator\_}.
  \item Test set usato una volta; importanza per permutazione sul test; \texttt{joblib.dump/load} per salvare.
\end{itemize}
\end{appunti}
